\documentclass[UTF8,11pt,a4paper]{ctexart}
\usepackage[a4paper,left=1.45cm,right=1.45cm,top=1.55cm,bottom=1.65cm]{geometry}
\usepackage{amsmath,amssymb,bm}
\usepackage{graphicx}
\usepackage{booktabs,tabularx,array,longtable}
\usepackage{enumitem}
\usepackage[dvipsnames,svgnames]{xcolor}
\usepackage{tikz}
\usetikzlibrary{arrows.meta,positioning,calc,fit,shapes.geometric}
\usepackage[most]{tcolorbox}
\usepackage{hyperref}
\usepackage{caption}
\usepackage{float}
\usepackage{fancyhdr}
\usepackage{titlesec}
\usepackage{microtype}
\IfFontExistsTF{Times New Roman}{\setmainfont{Times New Roman}}{\setmainfont{Liberation Serif}}
\IfFontExistsTF{Noto Sans CJK SC}{\setCJKmainfont{Noto Sans CJK SC}}{}
\IfFontExistsTF{Noto Sans CJK SC}{\setCJKsansfont{Noto Sans CJK SC}}{}
\IfFontExistsTF{DejaVu Sans Mono}{\setmonofont{DejaVu Sans Mono}}{\setmonofont{Latin Modern Mono}}
\definecolor{mainblue}{HTML}{0B5CAD}
\definecolor{softblue}{HTML}{EAF4FF}
\definecolor{darktext}{HTML}{1F2937}
\definecolor{softgray}{HTML}{F5F7FA}
\definecolor{accent}{HTML}{F59E0B}
\definecolor{goodgreen}{HTML}{2E7D32}
\definecolor{badred}{HTML}{B91C1C}
\definecolor{purple}{HTML}{6D5BD0}
\hypersetup{colorlinks=true,linkcolor=mainblue,urlcolor=DarkOrange,citecolor=mainblue}
\setlength{\parindent}{2em}
\setlength{\headheight}{14pt}
\setlength{\parskip}{0.22em}
\setlist[itemize]{leftmargin=2em,itemsep=0.23em,topsep=0.25em}
\setlist[enumerate]{leftmargin=2.2em,itemsep=0.23em,topsep=0.25em}
\renewcommand{\arraystretch}{1.2}
\captionsetup{font=small,labelfont=bf}
\titleformat{\section}{\Large\bfseries\sffamily\color{mainblue}}{\thesection}{0.7em}{}
\titleformat{\subsection}{\large\bfseries\sffamily\color{darktext}}{\thesubsection}{0.6em}{}
\titleformat{\subsubsection}{\normalsize\bfseries\sffamily\color{darktext}}{\thesubsubsection}{0.55em}{}
\pagestyle{fancy}
\fancyhf{}
\fancyhead[L]{\small\sffamily Back to Basics: Just Image Transformers}
\fancyhead[R]{\small\sffamily \leftmark}
\fancyfoot[C]{\small\thepage}
\renewcommand{\headrulewidth}{0.4pt}
\newtcolorbox{infobox}[1]{enhanced,breakable,colback=softblue,colframe=mainblue!70!black,boxrule=0.8pt,arc=2mm,left=1.5mm,right=1.5mm,top=1mm,bottom=1mm,title=\bfseries #1}
\newtcolorbox{keybox}[1]{enhanced,breakable,colback=softgray,colframe=darktext!45,boxrule=0.6pt,arc=2mm,left=1.5mm,right=1.5mm,top=1mm,bottom=1mm,title=\bfseries #1}
\newtcolorbox{notebox}[1]{enhanced,breakable,colback=accent!8,colframe=accent!80!black,boxrule=0.7pt,arc=2mm,left=1.5mm,right=1.5mm,top=1mm,bottom=1mm,title=\bfseries #1}
\newcommand{\term}[1]{\textbf{\textcolor{mainblue}{#1}}}
\begin{document}

\section{论文基本信息}
\begin{center}
\begin{tabularx}{0.94\linewidth}{>{\bfseries}p{0.18\linewidth}X}
\toprule
题目 & Back to Basics: Just Image Transformers\\
课件日期 & 2025-11-21\\
研究方向 & 扩散生成 / 纯图像 Transformer\\
一句话概括 & 重新审视生成模型是否必须依赖复杂压缩潜空间,以纯图像 Transformer 在像素或简单表示上学习生成分布。\\
\bottomrule
\end{tabularx}
\end{center}

\begin{infobox}{摘要要点}
重新审视生成模型是否必须依赖复杂压缩潜空间,以纯图像 Transformer 在像素或简单表示上学习生成分布。 本说明按“研究问题—核心机制—基础公式—实验阅读—局限与优化”的顺序重新组织，不保留原始材料的封面、目录和页码对应。
\end{infobox}

\section{研究问题与动机}
这项工作的出发点不是简单地替换某个网络层，而是识别现有范式中最影响\term{质量、效率、可靠性或可部署性}的瓶颈。结合论文所讨论的问题，可以将研究动机概括为以下几点：
\begin{itemize}
\item 对比 /v-pred ε 揭传统痛点,让读者能迅速get 到核心idea
\item JiT 确立 “ Diffusion + Transformer” 范式
\item 可推广至蛋白质、分子等难设计 Tokenizer 领域
\item JiT 无额外依赖,原生数据高效生成
\item “少即是多” 思路具 CV 领域价值
\item 凸显 “直接预测干净数据” 的价值
\end{itemize}

\begin{notebox}{理解重点}
阅读这类论文时，应把“问题是否真实存在”“方法是否直接解决该问题”“实验是否排除了其他解释”分开判断。仅凭更大的模型、更多数据或更长训练得到的提升，不能自动视为结构创新带来的收益。
\end{notebox}

\section{核心思想与整体流程}
重新审视生成模型是否必须依赖复杂压缩潜空间,以纯图像 Transformer 在像素或简单表示上学习生成分布。 从信息流角度看，其基本流程可以整理为下图。

\begin{figure}[H]
\centering
\begin{tikzpicture}[>=Latex,node distance=0.78cm]
\tikzstyle{procstep}=[draw,rounded corners,fill=softblue,align=center,minimum width=3.05cm,minimum height=0.95cm,text width=2.75cm]
\node[procstep] (a) {图像表示};
\node[procstep,right=of a] (b) {纯 Transformer};
\node[procstep,right=of b] (c) {噪声/流训练};
\node[procstep,right=of c] (d) {图像生成};
\draw[->,thick,mainblue] (a)--(b);
\draw[->,thick,mainblue] (b)--(c);
\draw[->,thick,mainblue] (c)--(d);
\end{tikzpicture}
\caption{核心信息流与处理阶段。}
\end{figure}

\subsection{方法组成与信息流}
\begin{itemize}
\item JiT - 极简架构设计(大patch Transformer )
\item 大 patch 保留局部语义完整性,减计算量,分辨率不影响长度
\item 无 Tokenizer 、预训练等,像素端到端训
\item 图像划非重叠大 patch ,降序列长度
\item 核心组件仅线性嵌入等,无额外模块
\item 对比 /v-pred ε 揭传统痛点,让读者能迅速get 到核心idea
\item JiT 确立 “ Diffusion + Transformer” 范式
\end{itemize}

\begin{keybox}{结构解析}
\begin{itemize}
\item 极简架构设计（大patch Transformer ） ； 图像划非重叠大 patch ，降序列长度 ； 256×256 用 16×16 patch （长 256 ） ； 512×512 用 32×32 patch （长 256 ） ； 核心组件仅线性嵌入等，
\end{itemize}
\end{keybox}

上述内容以文字方式保留方法结构，避免把整张幻灯片作为独立页面嵌入正文。

\section{数学与机制理解}
本节给出理解该工作的基础数学结构。公式用于解释方法所属范式，并不替代原论文中的完整推导。
\begin{equation}
x_t=\alpha_t x_0+\sigma_t\varepsilon,\qquad \varepsilon\sim\mathcal N(0,I)
\end{equation}
扩散类方法把真实样本与噪声按时间混合，模型再学习逆向去噪、速度场或条件流。不同论文的主要差异通常落在表示空间、网络骨干和训练目标上。

\begin{keybox}{从公式回到方法}
应重点观察论文究竟改变了公式中的哪一部分：是输入表示、连接矩阵、条件变量、参数化方式、训练目标，还是求解与调度过程。真正的创新通常可以在这一层被准确定位。
\end{keybox}


\section{实验结果应如何阅读}
\begin{itemize}
\item 多场景验证:单分辨率生成、跨分辨率生成、高维大patch 测试
\item FID-50K (生成质量)、计算量(Gflops )
\item 训练数据与设置:ImageNet 数据集
\item JiT-B/L/H/G 四种模型尺寸
\item 参数规模(M )、跨分辨率表现
\item 对比 /v-pred ε 揭传统痛点,让读者能迅速get 到核心idea
\item JiT 确立 “ Diffusion + Transformer” 范式
\end{itemize}

\begin{notebox}{结果解读}
\begin{itemize}
\item 实验验证 ； 训练数据与设置：ImageNet 数据集 ； （256×256/512×512/1024×1024 ） ； JiT-B/L/H/G 四种模型尺寸 ； 核心性能指标： ； FID-50K （生成质量）、计算量（Gflops ） ； 参数规模（M ）、跨……
\end{itemize}
\end{notebox}

该部分以文字概括实验结论与比较条件，避免用整页截图代替分析。
实验部分不应只看单一最好数字。还需要核对模型规模、训练数据、训练步数、采样或解码预算、硬件条件以及是否使用额外蒸馏、检索或后处理。只有比较条件接近时，结果才能真正说明方法本身的优势。

\section{优点、局限与可优化空间}
\subsection{主要优点}
\begin{itemize}
\item \term{问题与方法对应明确}：论文围绕一个可识别的核心瓶颈展开，方法结构能够直接映射到该瓶颈。
\item \term{可复用的设计思想}：即使不完整复现模型，其中的分解、稀疏化、递归、条件控制或系统优化思想也可迁移到相邻任务。
\item \term{实验提供了多角度证据}：实验通常同时展示主结果、效率比较、案例或消融，使结论不完全依赖单一指标。
\end{itemize}

\subsection{局限与进一步优化}
\begin{itemize}
\item 需要进一步区分\term{结构收益}与更大算力、更多数据、不同训练技巧带来的收益，并在等参数、等计算预算下进行严格比较。
\item 对超参数、输入分布和模型规模的敏感性仍应系统分析；在一个数据集上的最优配置不一定能直接迁移。
\item 实际部署还要考虑显存峰值、并发、延迟抖动、失败案例和鲁棒性，而不仅是离线平均指标。
\item 可尝试将该方法与蒸馏、量化、缓存复用、动态路由或更强的表示学习结合，验证不同优化是否互补。
\end{itemize}

\section{结论}
总体而言，重新审视生成模型是否必须依赖复杂压缩潜空间，以纯图像 Transformer 在像素或简单表示上学习生成分布。。 进一步需要注意：对比 /v-pred ε 揭传统痛点,让读者能迅速get 到核心idea;JiT 确立 “ Diffusion + Transformer” 范式。

\end{document}
